\chapter{The console, and scripting}
\label{ch:console}

Every button in MeerK40t is a front end for a console command. That fact is what
makes the program scriptable, reproducible and --- once you are comfortable ---
much faster to use. This chapter teaches the console as a tool rather than as a
reference; the complete command tables are in Appendix~\ref{app:console}.

\section{Opening the console and reading it}

\begin{itemize}
  \item Open the \panel{Console} pane with the ribbon's \btn{Console} button, or
        from \menu{Panes>Console}. (Many people bind \key{F12} to
        \cmd{window open Console} in the Keymap window; it is not bound by
default.)
  \item Type a command and press \key{Enter}. Output appears above the input line,
        coloured, with clickable links when a message contains a URL.
  \item Up and Down arrows walk your command history, which is kept between
        sessions --- the last fifty commands are stored in \file{cmdhistory.log}
        in the program's working directory.
  \item \btn{Tab} completes command names: type \cmd{sca} then \key{Tab} and the
        console offers the candidates.
  \item \cmd{cls} or \cmd{clear} empties the output; \cmd{console\_font} changes
        the typeface, and \cmd{console\_font -r} resets it.
\end{itemize}

\noindent The console is also the quiet witness of everything the program does.
When a job is planned, the stages go past in the console; when a device
disconnects, the reason is very often printed here first. If you report a problem,
the console text is the first thing anyone will ask for.

\begin{tipbox}[Record a session]
\cmd{channel save console} writes everything printed to the console channel into a
timestamped file --- the polite way to hand a support question to someone else.
\cmd{channel list} shows which channels exist, and \cmd{channel open <name>}
watches an individual one (device traffic, for example).
\end{tipbox}

\section{Asking the console for help}

\begin{lstlisting}[style=konsole]
help                 # every command, grouped by the kind of data it acts on
help rect            # one command, with its arguments and options
?                    # same as help
find rotary          # every command and option containing "rotary"
?? rotary            # same as find
\end{lstlisting}

\noindent \cmd{help} with no argument is the real index: commands are grouped by
category, and the groupings tell you how the console is structured --- some
commands create data, some act on selected elements, some act on a plan.

\section{The syntax rules that matter}

\begin{center}\small
\begin{tabular}{@{}L{4.0cm}L{10.2cm}@{}}
\toprule
\textbf{Rule} & \textbf{Example and effect} \\
\midrule
Units are written with the number & \cmd{rect 2cm 2cm 1cm 1cm} makes a
1\unit{cm} square 2\unit{cm} in from the origin. \cmd{mm}, \cmd{cm},
\cmd{inch}, \cmd{mil}, \cmd{deg} and percentages all work. \\
One command per line & A newline ends a command. Inside \cmd{-e "..."} on the
command line, a newline separates commands too. \\
\cmd{|} separates commands on one line & \cmd{rect 0 0 1cm 1cm | cut} runs both.
A \cmd{|} inside quotes is not a separator. \\
A leading \cmd{.} silences the echo & \cmd{.timer 1 0 spool jog} runs without
printing the command itself --- useful in scripts that would otherwise flood the
console. \\
Pipelines pass data along & A command that produces elements hands them to the
next command that accepts elements. \cmd{rect 0 0 1cm 1cm engrave -s 15} creates a
rectangle and immediately groups it into an engrave operation. \\
\cmd{; } belongs to aliases and stored batches & It is not a command separator in
ordinary use. \\
\bottomrule
\end{tabular}
\end{center}

\section{The job pipeline in one line}

The most useful sentence in MehrK40t's console is this:

\begin{lstlisting}[style=konsole]
planz clear copy preprocess validate blob preopt optimize spool
\end{lstlisting}

\noindent It builds a plan (\cmd{planz}), clears it, copies the enabled operations
in, preprocesses and validates them, converts them to cutcode (\cmd{blob}),
prepares and runs the optimisers, and sends the result to the spooler. It is
exactly what the \btn{Start} button does. There are two built-in aliases for
close relatives of it:

\begin{lstlisting}[style=konsole]
burn        # the chain above
simulate    # the same without spooling, then opens the Simulation window
\end{lstlisting}

\noindent Substitute \cmd{save_job <file>} for \cmd{spool} and the same chain
writes a device job file instead of running the machine:

\begin{lstlisting}[style=konsole]
plan0 clear copy preprocess validate blob preopt optimize save_job /tmp/tag.egv
\end{lstlisting}

\noindent \cmd{plan} with no name means the default plan, \cmd{plan0}; \cmd{planz}
addresses the plan called \cmd{z}. Naming plans is how you keep several jobs in
flight at once --- the Laser panel's \tabname{Plan} tab is a view of the same
objects.

\section{A complete example, typed}

\begin{lstlisting}[style=konsole]
rect 20mm 20mm 60mm 30mm          # the tag outline
circle 26mm 33mm 2mm              # a hanging hole (radius form)
select 0                          # nothing selected: element 0 is the rect
cut -s 8 -p 70                    # cut it: 8 mm/s, 70% power
select all
engrave -s 30 -p 25               # engrave everything else
planz clear copy preprocess validate blob preopt optimize spool
\end{lstlisting}

\noindent Note the shape of it: create, select, group into an operation with its
settings, then run the standard pipeline. Once you can read that, you can write
it.

\section{Settings from the console}

\cmd{set} is the single entry point for scalar settings. With no arguments it
lists them; with a key and value it assigns:

\begin{lstlisting}[style=konsole]
set                                # list the variables in the current context
set bedwidth 400mm                 # set a device value
set -p balor rotary_active True    # set it in the device's own context
set -p / rotary_diameter_mm 80     # explicit path form
flush                              # write settings to disk now
setting_export /home/me/backup     # copy the whole configuration out
setting_import /home/me/backup/meerk40t_2026-09-17_10_00_00.cfg
\end{lstlisting}

\noindent There is no general-purpose \cmd{get}: to read a value, list the
context with \cmd{set} and read the line you want. Device settings are covered in
Appendix~\ref{app:drivers}; the important habit is to use \cmd{set -p <device>}
rather than changing values globally.

\section{Selecting and acting from the console}

Element commands act on the current selection, and the console gives you several
ways to build one:

\begin{lstlisting}[style=konsole]
element            # the currently selected elements
element*           # all elements
element~           # everything that is not selected
element0,3,4       # specific elements by index
element* select    # select everything (select+ adds, select- removes, select^ toggles)
range 2 8          # indices 2 to 8
filter speed>=10   # operations matching an expression
tree selected delete
\end{lstlisting}

\noindent This is worth playing with rather than reading: create a few shapes and
watch the different selectors change what a following \cmd{translate 5mm 0}
affects.

\section{Aliases: your own commands}

An alias is a stored command line, and the console ships with some useful ones:

\begin{lstlisting}[style=konsole]
alias                      # list the aliases
burn                       # alias: the full plan-and-spool chain
simulate                   # alias: plan and open the simulation
alias myburn cut -s 6 -p 80 ; planz clear copy preprocess validate blob preopt optimize spool
\end{lstlisting}

\noindent Note that \cmd{;} separates the individual commands \emph{inside} an
alias body --- that is what the semicolon is for. Default aliases can be reset if
you break them (\menu{Help} or the bind/alias commands), and keyboard bindings
ultimately call these same console commands, which is why the Keymap window can
bind a key to anything you can type.

\section{Scripts: files, batches and the CLI}

There are three ways to run a list of commands without typing them, and they are
appropriate in different places.

\begin{description}[leftmargin=1.4em,style=nextline]
  \item[A command file] \cmd{execute /home/me/scripts/tag.txt} runs each line of
        the file as a console command. Lines beginning with \cmd{\#} are
        comments. This is the general-purpose scripting route.
  \item[The command line] \cmd{-e "command"} may be repeated, and \cmd{-b <file>}
        runs a file of commands at startup:
\begin{lstlisting}[style=konsole]
meerk40t -z -e "load /home/me/tag.svg" \
         -e "planz clear copy preprocess validate blob preopt optimize spool" -a -q
\end{lstlisting}
        That invocation loads a design, builds the job, starts the laser and quits
        when the spooler finishes --- headless, from a script or a scheduler.
  \item[The Auto-Exec window] \menu{Tools>Auto-Exec} (titled \emph{File startup
        commands}) stores commands inside the project file itself, and runs them
        when that file is loaded, if \btn{Execute on load} is ticked. The
        \btn{Execute commands} button runs them immediately. Its small \cmd{*}
        menu offers ready-made entries such as \emph{Start simulation},
        \emph{Burn content on current device} and \cmd{device activate} for each
        of your devices. This is the right place for the three lines a particular
        project always needs.
\end{description}

\noindent Useful meta-commands for scripts: \cmd{wait <seconds>} pauses,
\cmd{loop <from> <to> <command>} repeats, \cmd{threaded <command>} runs something
in the background, and \cmd{timer <times> <duration> <command>} repeats on a
schedule.

\section{Batch production from a spreadsheet}

The \btn{CSV Batch Run} window (\menu{Edit>Settings>CSV Batch Run}) is a
production tool: it reads a CSV file, exposes the columns as variables, and can
burn a job once per row.

\begin{walkthrough}{Numbered tags for a list of names}
\begin{enumerate}[label=\arabic*.,leftmargin=1.4em]
  \item Prepare a CSV with a header row, for example \cmd{name,serial}, and one
        row per item.
  \item Design the job with text elements that use those variables ---
        \cmd{{name}} and \cmd{{serial}} --- exactly as wordlists do
        (Chapter~\ref{ch:text}).
  \item Open the window and choose the \emph{File}. Press \btn{Import}.
  \item Set \emph{First row} to \emph{Auto}, \emph{Data} or \emph{Headers}
        depending on your file, and walk the rows with \btn{Prev}, \btn{Next} and
        \btn{Go to}; the header shows \emph{Row x of y}.
  \item Check the two preview panes: \emph{Preview (text with {variables})} shows
        what the design will say, and \emph{Current row values} shows what the
        variables currently hold.
  \item Press \btn{Apply row to design} to see one row in place without burning.
  \item Press \btn{Run current row} to burn just that row, or \btn{Run all rows}
        to work through the file --- with
        \emph{After each job completes, run the next row automatically} controlling
        whether it waits for you between jobs.
  \item \btn{Stop batch} stops the run; it does not unload the file, so you can
        resume where you left off.
\end{enumerate}
\end{walkthrough}

\noindent The window builds each job with the same plan chain discussed above
(with \cmd{threaded} in front when appropriate), which is why it behaves exactly
like pressing \btn{Start} --- including honouring your optimisation settings.

\section{Driving the interface from the console}

Two more command families are worth knowing because they turn repetitive setup
into one line:

\begin{lstlisting}[style=konsole]
window toggle Console            # open or close a window by name
window open Simulation 0 1 1     # open Simulation for plan 0 with flags
window reset *                   # forget all stored window positions
pane show console                # show a pane
pane hide tree ; pane float jog  # ... and the rest of the pane controls
page modify                      # switch the ribbon page
arm ; startjob                   # arm and start (same as the Ctrl+Alt+Shift+A binding)
\end{lstlisting}

\begin{tryit}{Script a repeatable production run}
\begin{enumerate}[label=\arabic*.,leftmargin=1.4em]
  \item Load the project you want to repeat and confirm it runs normally by
        pressing \btn{Start}.
  \item In the console, type \cmd{planz clear copy preprocess validate blob
        preopt optimize spool} and watch the stages go past. Nothing should
        differ from the button.
  \item Write those commands into a file called \file{run-tag.txt}, one per line,
        with a comment line at the top explaining what it is for.
  \item Test it with \cmd{execute /full/path/run-tag.txt}.
  \item Put the same commands into the Auto-Exec window and tick
        \btn{Execute on load}. Now loading the project prepares the job for you.
  \item Finally, try the headless form from a terminal: \cmd{meerk40t -z -e
        "load /full/path/tag.svg" -e "burn" -a -q}. This is the same job with no
        windows at all.
  \item Save the command file next to the project. In six months, it will be the
        only part of this exercise you still need.
\end{enumerate}
\end{tryit}
